% GENERATED FILE. Edit canonical lessons, then run npm run generate:books.
% canonical-source-hash: fnv1a64:3743da6331305aa6
% canonical-lessons: ZH-C418-zhiji, ZH-C418-shouxu, ZH-C418-tingfei, ZH-W418-ting-stop, ZH-C418-banli, ZH-C418-chaozhong

\chapter{Checking In}
\label{ch:zh-checking-in}
\hlchaptermodality{418}

\begin{chapteropening}
\textbf{By the end of this chapter:} \emph{I can check in for a flight, go through the formalities, say flights are grounded and a bag is too heavy.}

Complete the last lesson of chapter 418: I can check in for a flight, go through the formalities, say flights are grounded and a bag is too heavy.
\end{chapteropening}

\section[zhíjī]{zhíjī — to check in for a flight}

\label{lesson:ZH-C418-zhiji}



\begin{quote}
Before the new one: say the Mandarin for to declare goods at customs, then the Mandarin for duty-free.
\end{quote}

\subsection*{You'll want to know: zhíjī}
\emph{zhíjī} — \textquotedblleft{}to check in for a flight\textquotedblright{}.

\begin{quote}
\emph{wǒmen xiān zhíjī ránhòu qù ānjiǎn} — First we check in, then we go through security.
\end{quote}

Its characters come later; for now it is a word for the ear.

\begin{cousinweb}
\emph{zhíjī} is built from \emph{zhí} 'to be on duty' and \zh{机} 'a plane', both inherited from Middle Chinese, \emph{driH} and \emph{kj+j} (Guangyun, in Baxter's transcription), which glosses the first 'to hold' and 'to meet with'.
\end{cousinweb}

\begin{grammarlens}[title={What you've built}]
One more word for getting around.
\end{grammarlens}

\subsection*{Your turn}
\begin{itemize}
\raggedright
  \item \emph{Say it:} \emph{zhíjī}
  \item \emph{Say it:} \emph{zhíjī}, once more
  \item \emph{Say it:} \emph{wǒmen xiān zhíjī ránhòu qù ānjiǎn} — First we check in, then we go through security.
  \item \emph{Recall:} say the Mandarin for to declare goods at customs, then the Mandarin for duty-free, then say \emph{zhíjī} again
\end{itemize}

\begin{tcolorbox}[breakable,colback=teal!4,colframe=teal!35!black,title={Before you move on}]
What does zhíjī mean? (\textquotedblleft{}To check in for a flight\textquotedblright{}.) Say it once more.
\end{tcolorbox}
\section[shǒuxù]{shǒuxù — formalities; procedure}

\label{lesson:ZH-C418-shouxu}



\begin{quote}
Before the new one: say the Mandarin for duty-free, then the Mandarin for to check in for a flight.
\end{quote}

\subsection*{You'll want to know: shǒuxù}
\emph{shǒuxù} — \textquotedblleft{}formalities; procedure\textquotedblright{}.

\begin{quote}
\emph{chūguó de shǒuxù hěn máfan} — The formalities for going abroad are a real bother.
\end{quote}

Its characters come later; for now it is a word for the ear.

\begin{cousinweb}
\emph{shǒuxù} is built from \zh{手} 'a hand' and \emph{xù} 'to carry on', both inherited from Middle Chinese, \emph{syuwX} and \emph{zjowk} (Guangyun, in Baxter's transcription), which glosses the second 'to continue, to join'; how the pair came to mean formalities is not settled.
\end{cousinweb}

\begin{grammarlens}[title={What you've built}]
One more word for getting around.
\end{grammarlens}

\subsection*{Your turn}
\begin{itemize}
\raggedright
  \item \emph{Say it:} \emph{shǒuxù}
  \item \emph{Say it:} \emph{shǒuxù}, once more
  \item \emph{Say it:} \emph{chūguó de shǒuxù hěn máfan} — The formalities for going abroad are a real bother.
  \item \emph{Recall:} say the Mandarin for duty-free, then the Mandarin for to check in for a flight, then say \emph{shǒuxù} again
\end{itemize}

\begin{tcolorbox}[breakable,colback=teal!4,colframe=teal!35!black,title={Before you move on}]
What does shǒuxù mean? (\textquotedblleft{}Formalities; procedure\textquotedblright{}.) Say it once more.
\end{tcolorbox}
\section[\texorpdfstring{\zh{停飞} (tíngfēi)}{停飞 (tíngfēi)}]{\zh{停飞} (tíngfēi) — to be grounded (of planes)}

\label{lesson:ZH-C418-tingfei}



\begin{quote}
Before the new one: say the Mandarin for to check in for a flight, then the Mandarin for the formalities.
\end{quote}

\subsection*{You'll want to know: \zh{停飞}}
\textbf{\zh{停飞}} — \emph{tíngfēi} — \textquotedblleft{}to be grounded (of planes)\textquotedblright{}.

\begin{quote}
\emph{fēijī tíngfēi le wǒmen zuò gāotiě qù Běijīng} — The planes are grounded: we're taking the high-speed train to Beijing.
\end{quote}

The new character, in pieces: \zh{停} — eleven strokes in two pieces: on the left the person piece \zh{亻}; on the right a dot and a bar, a small \zh{口}, a short stroke down and a bar that hooks like a lid, a bar, and a vertical with a hook. The next lesson puts it in your hand.

\begin{cousinweb}
\textbf{\zh{停飞}} (\emph{tíngfēi}) is built from \zh{停} 'to stop' and \zh{飞} 'to fly', both inherited from Middle Chinese, \emph{deng} and \emph{pj+j} (Guangyun, in Baxter's transcription), which glosses the first 'to rest, to stop' and the second 'to fly, to soar'.
\end{cousinweb}

\begin{grammarlens}[title={What you've built}]
One more word for getting around.
\end{grammarlens}

\subsection*{Your turn}
\begin{itemize}
\raggedright
  \item \emph{Say it:} \emph{tíngfēi}
  \item \emph{Say it:} \emph{tíngfēi}, once more
  \item \emph{Say it:} \emph{fēijī tíngfēi le wǒmen zuò gāotiě qù Běijīng} — The planes are grounded: we're taking the high-speed train to Beijing.
  \item \emph{Recall:} say the Mandarin for to check in for a flight, then the Mandarin for the formalities, then say \emph{tíngfēi} again
\end{itemize}

\begin{tcolorbox}[breakable,colback=teal!4,colframe=teal!35!black,title={Before you move on}]
What does \zh{停飞} mean? (\textquotedblleft{}To be grounded (of planes)\textquotedblright{}.) Say it once more.
\end{tcolorbox}
\section[\texorpdfstring{\zh{停} (tíng)}{停 (tíng)}]{\zh{停} — the person piece, then a tall right side}

\label{lesson:ZH-W418-ting-stop}



\begin{quote}
The word first, then the character inside it.

\begin{itemize}
\raggedright
  \item \emph{Recall:} say \emph{tíngfēi} — \textbf{R1}, one lesson back, before it decays
\end{itemize}

\begin{quote}
\zh{你} \zh{高}
\end{quote}

\emph{Read it:} each one aloud, with nothing said first — \emph{nǐ}, \emph{gāo}
\end{quote}

\subsection*{Script — the person piece, then a tall right side}
\hlblockfigure{\detokenize{figures/ZH-W418-ting-stop-filmstrip.pdf}}{How \zh{停} is written, stroke by stroke}

\begin{quote}
\zh{停}
\end{quote}

\textbf{Two pieces, left before right.} Count the pieces, not the strokes.

\textbf{The left piece, \zh{亻}, two strokes.} The same two as in \zh{你}.

\textbf{The right piece, nine strokes.} A dot and a bar, then a small \zh{口}, three strokes: the top of \zh{高}. Under them a short stroke down and a bar that turns down and hooks: a lid. Then a bar across, and a vertical down the middle that hooks to the left at the foot.

2 + 9 = eleven.

\subsection*{Your turn}
\begin{itemize}
\raggedright
  \item \emph{Write it:} \zh{停}, all eleven, with the model in front of you
  \item \emph{Cover:} write it again from memory, and count the pieces, not the strokes
\end{itemize}

\begin{tcolorbox}[breakable,colback=teal!4,colframe=teal!35!black,title={Before you move on}]
How many pieces does \zh{停} have? (\textbf{Two}.) Which comes first? (\textbf{\zh{亻}, the person piece}.) How does the last stroke end? (\textbf{With a hook to the left}.)

Source: \href{https://raw.githubusercontent.com/chanind/hanzi-writer-data/68d10a4b21150cae5e1ebbd223eed289cf32d90c/data/\%E5\%81\%9C.json}{Hanzi Writer Data \zh{停}, pinned stroke paths and medians}.
\end{tcolorbox}
\section[bànlǐ]{bànlǐ — to deal with the paperwork; to go through a procedure}

\label{lesson:ZH-C418-banli}



\begin{quote}
Before the new one: say the Mandarin for the formalities, then the Mandarin for to be grounded.
\end{quote}

\subsection*{You'll want to know: bànlǐ}
\emph{bànlǐ} — \textquotedblleft{}to deal with the paperwork; to go through a procedure\textquotedblright{}.

\begin{quote}
\emph{wǒmen qù nà jiā bīnguǎn bànlǐ rùzhù} — We're going to that hotel to check in.
\end{quote}

Its characters come later; for now it is a word for the ear.

\begin{cousinweb}
\emph{bànlǐ} is built from \emph{bàn} 'to see to' and \emph{lǐ} 'to put in order', both inherited from Middle Chinese, \emph{beanH} and \emph{liX} (Guangyun, in Baxter's transcription), which glosses the first 'to prepare' and the second 'to manage'.
\end{cousinweb}

\begin{grammarlens}[title={What you've built}]
One more word for getting around.
\end{grammarlens}

\subsection*{Your turn}
\begin{itemize}
\raggedright
  \item \emph{Say it:} \emph{bànlǐ}
  \item \emph{Say it:} \emph{bànlǐ}, once more
  \item \emph{Say it:} \emph{wǒmen qù nà jiā bīnguǎn bànlǐ rùzhù} — We're going to that hotel to check in.
  \item \emph{Recall:} say the Mandarin for the formalities, then the Mandarin for to be grounded, then say \emph{bànlǐ} again
\end{itemize}

\begin{tcolorbox}[breakable,colback=teal!4,colframe=teal!35!black,title={Before you move on}]
What does bànlǐ mean? (\textquotedblleft{}To deal with the paperwork; to go through a procedure\textquotedblright{}.) Say it once more.
\end{tcolorbox}
\section[chāozhòng]{chāozhòng — overweight; over the weight limit}

\label{lesson:ZH-C418-chaozhong}



\begin{quote}
Before the new one: say the Mandarin for to be grounded, then the Mandarin for to deal with the paperwork.
\end{quote}

\subsection*{You'll want to know: chāozhòng}
\emph{chāozhòng} — \textquotedblleft{}overweight; over the weight limit\textquotedblright{}.

\begin{quote}
\emph{nǐ de xiāngzi chāozhòng le} — Your suitcase is overweight.
\end{quote}

Its characters come later; for now it is a word for the ear.

\begin{cousinweb}
\emph{chāozhòng} is built from \emph{chāo} 'to go beyond' and \zh{重} 'heavy', both inherited from Middle Chinese, \emph{trhjew} and \emph{drjowngX} (Guangyun, in Baxter's transcription), which glosses the first 'to leap'.
\end{cousinweb}

\begin{grammarlens}[title={What you've built}]
One more word for getting around.
\end{grammarlens}

\subsection*{Your turn}
\begin{itemize}
\raggedright
  \item \emph{Say it:} \emph{chāozhòng}
  \item \emph{Say it:} \emph{chāozhòng}, once more
  \item \emph{Say it:} \emph{nǐ de xiāngzi chāozhòng le} — Your suitcase is overweight.
  \item \emph{Recall:} say the Mandarin for to be grounded, then the Mandarin for to deal with the paperwork, then say \emph{chāozhòng} again
\end{itemize}

\begin{tcolorbox}[breakable,colback=teal!4,colframe=teal!35!black,title={Before you move on}]
What does chāozhòng mean? (\textquotedblleft{}Overweight; over the weight limit\textquotedblright{}.) Say it once more.
\end{tcolorbox}
